\section{The Poisson experiment and fixed-sample risk}\label{sec:risk}

Choose a dyadic \(h\asymp(n\log n)^{-1/(s+d)}\), and set
\(\tau=h^s\), \(v=(3h)^d\), so that \(v\tau\asymp(n\log n)^{-1}\).
Write \(a_s=s/(s+d)\) and fix \(\eta,\gamma>0\) with
\(\gamma+\eta<1-2a_s\). Choose \(J\) so that
\(\epsilon=(3/4)^J/2<\min(\omega,1/(s+d))\).

Start with a \(\Pois(n/2)\) sample and give each point an
independent uniform label among six groups and an independent
continuous order mark. The six groups are used for selection,
low-mass correction, shape estimation, estimation of the
enlarged-cell mass, estimation of the evaluation-cell mass,
and estimation of the log density scale. Each group has
mean \(m=n/12\).
A local group count is computed on \(J_I^{\rm phys}\);
apply the inverse of \(T_I\) only after this restriction.
The local count is Poisson of mean \(mP_I\) and, conditional
on that count, the positions in unit coordinates are iid
of density \(w_I\). Ordering by independent marks applies a
permutation independent of the positions. These laws follow
by splitting the original count according to its point
categories. For disjoint physical supports, the generating
function of the counts for labels and support indicators is
\(\prod_j\exp(\lambda_j(z_j-1))\) with the appropriate means
\(\lambda_j\), by averaging the conditional multinomial
generating function. Given the counts, the marked positions
are iid within each category and independent across categories.
Thus the entire marked restrictions on disjoint physical
supports are independent. All cell estimators below use their
own marked restriction only.

Use the selector group to choose the low rule when its
count in \(J_I^{\rm phys}\) is less than \(D_{\rm sel}\log n\).
The low rule is valid at \(P_I\le A=C_A\log n/n\), and the
other rule at \(P_I\ge C_H\log n/n\).
First choose \(C_H\) sufficiently large for the nonlow rule.
Then take \(D_{\rm sel}\) so that for
\(P_I<C_H\log n/n\) the selector mean is at most half the
threshold, and \(C_A\) so that for \(P_I>A\) that mean
is at least twice the threshold. Poisson tails give failure
probability \(O(n^{-M})\) for any required fixed \(M\) by
increasing these constants. The validity ranges overlap.
In the nonlow rule, the shape group count \(N\) obeys
\(N\asymp nP_I\) and \(C_*\log n\le N\le n\) on an
event of probability at least \(1-O(n^{-M})\) in that
range, by the same concentration and \(P_I\le1\).
Outside the observed window use zero for the shape statistic.

\subsection{Low mass}
The local Poisson intensity in unit coordinates is
\(m w(x)\dd x\), where \(w=v f\circ T_I=P_I w_I\)
for a cell of positive mass \(P_I\). Its projection and residual
are \(q=P_I q_I\) and \(e=P_I E_I\).
For a point from the full correction label, the position
on the torus has density \(f\). Map it to \(T_I^{-1}\) of
its location if in \(J_I^{\rm phys}\), and to the outside
point of Section~\ref{sec:angular} otherwise. This gives
precisely the parent observation probability and restriction
used there. Put \(r_*=A/K^2\). By
\(\psi(P_I q_I)=P_I\psi(q_I)+q\log P_I\) for \(P_I>0\),
\begin{equation}\label{eq:low-decomp}
 F_I=\int_W [\psi(q)-q\log A]
       +\left(\mu_I-\int_W q\right)\log(P_I/A)
       +\mu_I\log(A/v).
\end{equation}
At \(P_I=0\) all terms are interpreted as zero.
The first term is estimated by Section~\ref{sec:angular},
with residual \(P_I E_I\) and indicator
\(\one_{\{\inf_W |q|<P_I/2\}}=\nu_I\), also valid at \(P_I=0\)
under the zero-mass convention.
The final term is estimated linearly with the score \(\one_W\).

For the middle term, a degree-\(K\) polynomial \(P_K\)
approximates \(x\log x\) uniformly on \([0,1]\) to error
\(C/K^2\), with coefficient norm at most \(C^K\).
For a construction, put \(x=(1+\cos\theta)/2\).
The cosine coefficients of \(\log x\) at modes \(j\ge1\) are
\(2(-1)^{j-1}/j\), by the radial
logarithm expansion of \(\log|1+\rho e^{i\theta}|^2-\log 4\)
as \(\rho\) increases to one. The series converges in \(L^2\)
by square summability, and its pointwise limit off
\(\theta=\pi\) identifies the \(L^2\) limit. Multiplication by
\(x\) combines two adjacent coefficients and the central one,
giving \(O(j^{-3})\) for \(j\ge2\). The absolutely summable
series for \(x\log x\) therefore has tail \(O(K^{-2})\)
uniformly, and its truncation has the coefficient bound by
the Chebyshev recurrence. Now \(B_{K-1}(x)=(P_K(x)-P_K(0))/x\)
is a polynomial with
\(x|B_{K-1}(x)-\log x|\le C/K^2\) for \(x>0\).
The multiplier \(\mu_I-\int_W q\) is the mean of
\[
 b(y)=\one_W(y)-\int_W K_\Omega(x,y)\dd x
\]
on the parent law, extended by zero outside \(\Omega\).
Its absolute value is \(O(P_I)\).
For clarity, the score \(b\) is represented in the statistic
as \(2D_b g_b-D_b\one_\Omega\), where
\(D_b\ge\max(1,\|b\|_\infty)\) and
\(g_b=\one_\Omega(1+b/D_b)/2\) belongs pointwise to \([0,1]\).
Use the factorial polynomial estimator for the product of
this mean and \(B_{K-1}(P_I/A)\).
Together with the linear term, its coefficient norm in the
scaled score means is \(O(C^K A\log n)\).
The bias from the replacement of the logarithm is at most
\(C A/K^2\), including zero mass by continuity of the
mass-multiplied target and the polynomial formula.
All terms have the parent denominator \(m^k\).
The factorial bound and Minkowski give a combined second
moment at most \(C n^\eta A^2\log^2 n\); decrease \(c_K\) if
needed for this combined bound. The bias is at most
\begin{equation}\label{eq:low-bias}
 C\bigl[P_I E_I+r_*(1+(P_I/r_*)^\epsilon
                          \nu_I)\bigr].
\end{equation}

\subsection{Nonlow mass}
Conditional on \(N\) in the observed range, the ordered
shape sample supplies the independent fractions in
Section~\ref{sec:angular} for \(N<n^\gamma\) and
Section~\ref{sec:high} for \(N\ge n^\gamma\).
The shape estimate of \(\int_W\psi(q_I)\) is multiplied
by the independent shape-mass count divided by \(m\), an
unbiased estimate \(\widehat P_I\) with variance \(P_I/m\).
On the good counts and good actual pilots, the resulting
conditional bias for \(P_I\int_W\psi(q_I)\) is bounded by
\begin{equation}\label{eq:mass-shape-bias}
 C\left[P_I E_I+\frac{P_I}{N\log n}
       \bigl(1+(N\log n)^\epsilon\nu_I\bigr)\right].
\end{equation}
Here \(r_*=\log n/(NK^2)\asymp (N\log n)^{-1}\).
Its conditional variance at these histories is
\(C n^\eta\log^C n\, P_I^2\) for \(N<n^\gamma\), and
\(C n^\eta\log^C n\,P_I/n\) for \(N\ge n^\gamma\).
Indeed \(\E\widehat P_I^2=O(P_I^2)\) on the valid
range. The product variance for an independent mass factor
and shape factor is this second moment times the shape
variance, plus the mass variance times the squared shape
mean. In the high range use \(N\asymp nP_I\) and the
\(C\log n\) mean bound from Section~\ref{sec:high}.
In the intermediate range the shape second moment alone
and \(nP_I\gtrsim\log n\) suffice.

Use two more independent groups for \(\mu_I\log b_I\).
Estimate \(\mu_I\) by the evaluation-cell count divided by
\(m\), and use the enlarged-cell count \(C_I\) in the
other group to form
\[
 L_{C_I}=\sum_{j=1}^{C_I}j^{-1}-\gamma_{\rm E}-\log(mv),
\]
where \(\gamma_{\rm E}\) is Euler's constant and the empty
sum is zero. For \(\lambda=mP_I\), the mean of the harmonic
sum equals \(\int_0^\lambda(1-e^{-u})\dd u/u\), by using
\(\sum_{j=1}^k j^{-1}=\int_0^1 (1-t^k)\dd t/(1-t)\)
and the Poisson generating function. For \(C_I\ge\lambda/2\) and \(\lambda\ge2\),
the harmonic sum differs from the sum at \(\lfloor\lambda\rfloor\)
by at most \(C|C_I-\lfloor\lambda\rfloor|/\lambda\).
The complementary event has exponentially small probability
and squared difference at most \(C(1+\log\lambda)^2\).
The expected squared difference is therefore at most \(C/\lambda\).
This bounds the variance and shows that the mean harmonic sum
minus \(\log\lambda\) tends to \(\gamma_{\rm E}\) as
\(\lambda\to\infty\). Subtracting \(\log\lambda\) in its
integral representation and integrating the derivative
\(-e^{-\lambda}/\lambda\) to infinity gives
\(\gamma_{\rm E}+\int_\lambda^\infty e^{-u}\dd u/u\).
As \(\lambda/(mv)=b_I\), the bias of \(L_{C_I}\) for \(\log b_I\)
is \(\int_\lambda^\infty e^{-u}\dd u/u\).
The bias can be made \(O(n^{-M})\) on the nonlow range
by \(C_H\) large. Independence, \(\mu_I\le P_I\) and
\(|\log b_I|\le C\log n\) there give variance at most
\(C P_I\log^2 n/n\) for the product with the estimate of
\(\mu_I\).

\subsection{Summation and return to fixed sample size}

Let \(Z_I\) be the output selected in cell \(I\), clipped
to \([-n^{B_0},n^{B_0}]\) for a fixed sufficiently large
\(B_0\). Let \(\mathcal H_I\) be the sigma-field of the
local selector and, on nonlow routes, the count \(N\) and
actual shape pilots used by the rule. The good event consists
of a valid route and, for a nonlow rule, count comparability
and validity of the real pilots. It is \(\mathcal H_I\)-measurable
for fixed \(f\), with failure probability \(O(n^{-M})\).
The estimator only checks observed conditions. In particular
it uses the actual interval-pilot regional count for the
gate in Section~\ref{sec:high}, and the auxiliary path
probabilities there and in Section~\ref{sec:angular} are
built into the estimator.

The good-history second moments of the unclipped outputs
are bounded by a fixed power of \(n\) uniformly over \(f,I\).
This follows from the bounds on their biases, variance and
means, with \(|F_I|\le C\log n\) even at small mass since
\(P_I|\log P_I|\) is bounded. One may choose the clipping
exponent before making the failure probability arbitrarily
small: the powers in the good-history bounds are controlled
by the fixed \(\eta,\gamma,s,d\), with arbitrary concentration
constants absorbed by enlarging the sample-size cutoff.
Clipping changes each good-history mean by \(o(v\tau)\)
using the second moment divided by the clipping threshold,
and contracts the conditional variance by the independent-copy
formula. On bad histories the squared conditional bias and
the conditional variance are bounded by
\(C(n^{B_0}+\log n)^2\). Choose \(M>2B_0+3\).
The sum of the bad-history conditional bias contributions
in \(L^2\) is at most
\(C h^{-d}(n^{B_0}+\log n)n^{-M/2}=o(\tau)\).
For the variance sum the bound is
\(C h^{-d}n^{2B_0-M}=o(\tau^2)\).

Combining \eqref{eq:low-bias} and
\eqref{eq:mass-shape-bias}, on all good histories
\begin{equation}\label{eq:per-cell}
 \left|\E[Z_I\mid\mathcal H_I]-F_I\right|
 \le C P_I E_I+ C v\tau\left(1+(b_I/\tau)^\epsilon
                                 \nu_I\right).
\end{equation}
On low routes \(A/K^2\asymp v\tau\) and
\(P_I/(A/K^2)\asymp b_I/\tau\) when \(P_I>0\).
On the good nonlow histories \(N\asymp nP_I\) gives
\(P_I/(N\log n)\asymp v\tau\) and \(N\log n\asymp b_I/\tau\).
The residuals sum to \(O(\tau)\) by \eqref{eq:projection}.
The terms for \(b_I\le\tau\) and the constant term
also sum to \(O(\tau)\). For the remaining term partition
the amplitudes into disjoint dyadic bands
\(B/2<b_I\le B\) with \(B=2^k\tau\), \(k\ge1\).
At \(\xi=B/\tau\), Lemma~\ref{lem:bands} gives
a contribution at most
\[
 C\tau\,\xi^\epsilon(\xi^{-\omega}+h)
       +C \xi^{\epsilon-1}\sum_{B/2<b_I\le B}\sqrt v\,e_I .
\]
The largest \(\xi\) is at most \(C h^{-(s+d)}\)
since \(P_I\le1\). The dyadic sum of \(\xi^{\epsilon-\omega}\)
is bounded since \(\epsilon<\omega\), and that of
\(h\xi^\epsilon\) is \(O(h^{1-(s+d)\epsilon})\). The residuals remain in disjoint
bands and \(\xi^{\epsilon-1}\le1\). Therefore the sum of the
conditional biases has \(L^2\) norm \(O(\tau)\) by Minkowski,
including the bad-history contribution.

For the centered part, color the intersection graph of the
physical supports by at most \(D_{\rm col}=7^d\) colors,
using the degree bound from Section~\ref{sec:geometry}.
Within each color the local marked restrictions are independent.
Both \(Z_I\) and its conditional mean given \(\mathcal H_I\)
are measurable with respect to this local restriction.
Consequently the centered terms have mean zero and are
independent inside each color. By Cauchy--Schwarz across colors,
\[
 \E\left|\sum_I (Z_I-\E[Z_I\mid\mathcal H_I])\right|^2
 \le D_{\rm col}\sum_I \E\Var(Z_I\mid\mathcal H_I).
\]
On good low histories the sum of the conditional variances is
bounded by \(C h^{-d} n^\eta A^2\log^2 n\). Good intermediate
histories occur only for \(P_I\le C n^{-1+\gamma}\),
by \(N<n^\gamma\) and count comparability there.
The good intermediate variances thus sum, in expectation,
to at most \(C n^{-1+\gamma+\eta}\log^C n\),
using \(\sum_I P_I\le5^d\).
The good high histories contribute
\(C n^{-1+\eta}\log^C n\) by the same mass sum.
Together these bounds give
\[
 C\log^C n\bigl(h^{-d}n^\eta A^2+n^{-1+\gamma+\eta}\bigr)
   +o(\tau^2)=O(\tau^2).
\]
For the first term the strict power slack is
\(\eta<1-a_s\), and for the others it is
\(\eta+\gamma<1-2a_s\). Both follow from our choices.

The Poisson output \(\sum_I Z_I\) is bounded for each \(n\)
and measurable. The centered variance and conditional bias
bounds and \eqref{eq:transfer} give its root-mean-square
error at most \(C\tau\) for \(\int f\log f\), uniformly
over the class. The mesh cutoff \(h\le1/8\) and all
stochastic cutoffs are achieved by increasing a single
\(n_0(d,s,p,L)\), since all estimates above are uniform.

To use exactly \(n\) observations, draw \(M\sim\Pois(n/2)\)
independently. If \(M\le n\), apply the Poisson procedure to
the first \(M\) observations with independent auxiliary labels
and marks, and reverse its sign. Otherwise output zero.
Couple to an infinite iid sequence to see that the squared
error on \(M\le n\) is at most the squared error in the
complete Poisson experiment. On \(M>n\) the squared error
is \(H(f)^2\), bounded uniformly by \(\|f\|_2\le L\) and
\(|u\log u|\le C(1+u^2)\). The overflow probability is at most
\(\exp[-n(\log 2-1/2)]\) by the exponential moment at
\(\log 2\). The auxiliary-randomness average at each sample
vector is a finite measurable estimator. Jensen's inequality
gives no larger squared error for this deterministic estimator.
This proves the upper bound for \(2<s<d\). The likelihood
mixture of Section~\ref{sec:large-s} gives the rest of the
upper bound, completing Theorem~\ref{thm:main}.
